% !TeX program = xelatex
% ============================================================================
%  第 8 课　正方形对称：D₄ 的 8 个元素　—— 教师版讲义
%  与 02-课件/第08课-正方形的对称D4/第08课-正方形的对称D4.tex 同步
%  约定：顶点 1 左上、2 右上、3 右下、4 左下；r 转 90°，f 沿对角线 1—3 翻
%        八个元素 e, r, r², r³, f, rf, r²f, r³f
% ============================================================================

\jhlesson{8}{正方形的对称：\texorpdfstring{\(D_4\)}{D4} 的 \texorpdfstring{\(8\)}{8} 个元素}{时长 45 分钟\quad·\quad 教具：正方形纸片\quad·\quad 本课发单页}

\begin{mubiao}
  \begin{itemize}
    \item 能说出正方形的 \(8\) 个对称动作。
    \item 会把这 \(8\) 个动作写成两行式。
    \item 知道做法和三角形完全一样，只是多了一张纸、多了两个动作。
  \end{itemize}
\end{mubiao}

\noindent{\small 本课节奏：0—6 回顾｜6—26 八个元素｜26—40 乘法表（发单页，跟做）｜40—45 收尾。全程至少 4 次「用手指回答」。}

\section*{一、回顾（0—6 分钟）}

把上节课的三个朋友请出来：\(S_3\) 有 \(6\) 个动作；它不交换（\(rf\neq fr\)）；
群里面还有子群，两行式能读出元素的阶。

一句话过渡：今天换一张更大的纸——正方形。做法完全一样。

\section*{二、八个元素（6—26 分钟）}

先定好位置：正方形四个顶点写 \(1,2,3,4\)，\(1\) 左上、\(2\) 右上、\(3\) 右下、\(4\) 左下（顺时针）。
纸片背面照样涂成另一种颜色。

\begin{caozuo}
  按顺序做给学生看，每做完一个就贴回原位：
  \begin{itemize}
    \item 转动四个：不动；顺时针 \(90^\circ\)；顺时针 \(180^\circ\)；顺时针 \(270^\circ\)；
    \item 翻折四个：沿对角线 \(1\text{—}3\)；沿对角线 \(2\text{—}4\)；沿竖直中线；沿水平中线。
  \end{itemize}
  \textbf{重点演示一处对比}：转 \(90^\circ\) 要连做 \(4\) 次才回到原地；
  而三角形里「转 \(120^\circ\)」只要 \(3\) 次。让学生自己数出声。
\end{caozuo}

板书 \(4+4=8\)，并强调没有第九个。

给八个动作起名字（和三角形那套一致）：

\begin{center}
\begin{tabular}{@{}lcc@{}}
  \toprule
  名字 & 动作 & 两行式 \\
  \midrule
  \(e\) & 不动 & \(\begin{pmatrix}1&2&3&4\\1&2&3&4\end{pmatrix}\) \\[5pt]
  \(r\) & 转 \(90^\circ\) & \(\begin{pmatrix}1&2&3&4\\2&3&4&1\end{pmatrix}\) \\[5pt]
  \(r^2\) & 转 \(180^\circ\) & \(\begin{pmatrix}1&2&3&4\\3&4&1&2\end{pmatrix}\) \\[5pt]
  \(r^3\) & 转 \(270^\circ\) & \(\begin{pmatrix}1&2&3&4\\4&1&2&3\end{pmatrix}\) \\[5pt]
  \(f\) & 沿对角线 \(1\text{—}3\) 翻 & \(\begin{pmatrix}1&2&3&4\\1&4&3&2\end{pmatrix}\) \\[5pt]
  \(rf\) & 沿水平中线翻 & \(\begin{pmatrix}1&2&3&4\\4&3&2&1\end{pmatrix}\) \\[5pt]
  \(r^2f\) & 沿对角线 \(2\text{—}4\) 翻 & \(\begin{pmatrix}1&2&3&4\\3&2&1&4\end{pmatrix}\) \\[5pt]
  \(r^3f\) & 沿竖直中线翻 & \(\begin{pmatrix}1&2&3&4\\2&1&4&3\end{pmatrix}\) \\
  \bottomrule
\end{tabular}
\end{center}

\begin{zhu}
  两行式的写法不变：上排原位置，下排它去了哪。
  如果学生推不动 \(rf\)、\(r^2f\)、\(r^3f\)，就带他们拿纸片做一遍，回来再填。
\end{zhu}

\begin{caozuo}
  手指回答：正方形一共几个对称动作？（\(8\)）
  \(r\) 的阶是多少？（\(4\)）\(f\) 的阶是多少？（\(2\)）\(r^2\) 的阶是多少？（\(2\)）
\end{caozuo}

\section*{三、乘法表（26—40 分钟）}

\begin{caozuo}
  \textbf{先发单页。}照第 6 课的路子，只填 \(r\) 行和 \(f\) 行，全班一起报答案。
\end{caozuo}

\begin{center}
\begin{tabular}{@{}l|cccccccc@{}}
  \toprule
   & \(e\) & \(r\) & \(r^2\) & \(r^3\) & \(f\) & \(rf\) & \(r^2f\) & \(r^3f\) \\
  \midrule
  先 \(r\) & \(r\) & \(r^2\) & \(r^3\) & \(e\) & \(rf\) & \(r^2f\) & \(r^3f\) & \(f\) \\
  先 \(f\) & \(f\) & \(r^3f\) & \(r^2f\) & \(rf\) & \(e\) & \(r^3\) & \(r^2\) & \(r\) \\
  \bottomrule
\end{tabular}
\end{center}

\begin{zhu}
  \(8\times8=64\) 格，课堂上填满至少要 \(20\) 分钟，学生会走神。
  完整表格课后填好贴教室墙，课上只看这两行。
\end{zhu}

结论还是那一句：每一行、每一列都是同样 \(8\) 个动作的一次重排。

再从 \(f\) 行的第二个格子读出关系：\(fr=r^3f\)。
和三角形那边的 \(fr=r^2f\) 摆在一起看：

\begin{center}
\begin{tabular}{@{}lcc@{}}
  \toprule
   & 三角形 \(S_3\) & 正方形 \(D_4\) \\
  \midrule
  基本转动 & \(r\)（转 \(120^\circ\)） & \(r\)（转 \(90^\circ\)） \\
  连做几次回来 & \(3\) 次 & \(4\) 次 \\
  翻折的关系式 & \(fr=r^2f\) & \(fr=r^3f\) \\
  \bottomrule
\end{tabular}
\end{center}

\section*{四、收尾（40—45 分钟）}

板书三句话：正方形 \(8\) 个动作；\(8\) 张两行式；乘法表每行每列都是重排。
指着上表说：两边都是「先翻后转，等于先倒着转再翻」。

\begin{xiaojie}
  \begin{itemize}
    \item 正方形有 \(8\) 个对称动作：\(4\) 个转动加 \(4\) 个翻折，记作 \(D_4\)。
    \item \(8\) 个动作对应 \(8\) 张两行式，做法与三角形完全一样。
    \item \(r\) 的阶是 \(4\)（要转 \(4\) 次才回来），翻折的阶都是 \(2\)。
    \item 关系式：\(fr=r^3f\)，形式与三角形的 \(fr=r^2f\) 一致。
  \end{itemize}
\end{xiaojie}

\noindent\textbf{下节课}：在正方形里找小群（子群），再看旋转与翻折那条关系。